\documentclass[11pt]{article}
\usepackage[a4paper,margin=2.05cm]{geometry}
\usepackage[T1]{fontenc}
\usepackage{lmodern}
\usepackage{amsmath,amssymb}
\usepackage{xcolor}
\usepackage{microtype}
\usepackage[hidelinks]{hyperref}
\setlength{\emergencystretch}{3em}

\definecolor{ink}{HTML}{172033}
\definecolor{accent}{HTML}{385B7A}
\definecolor{soft}{HTML}{EEF5FA}
\color{ink}

\title{Complex phase retrieval in dimension four needs exactly eleven measurements}
\author{Literature answer to the open minimum in arXiv:1605.08034}
\date{11 August 2026}

\begin{document}
\maketitle

\begin{center}
\fcolorbox{accent}{soft}{\parbox{0.91\textwidth}{
\textbf{Status: literature-already-answered, exact value.}
No ten measurement vectors can perform standard phase retrieval in $\mathbb C^4$, while an explicit eleven-vector frame does.  Therefore
\[
N_{\min}(\mathbb C^4)=11.
\]}}
\end{center}

\section*{Source problem}
For vectors $a_1,\dots,a_N\in\mathbb C^4$, standard phase retrieval asks that
\[
 x\longmapsto \bigl(|a_1^*x|^2,\dots,|a_N^*x|^2\bigr)
\]
be injective modulo multiplication of $x$ by a scalar in the unit circle.  Wang--Xu state on source PDF page 5 that the smallest possible $N$ was unknown even for $d=4$ \cite{WX}.  Vinzant had already constructed an injective eleven-vector frame, whereas the best lower bound allowed ten.

\section*{Explicit later solution}
Huang's Theorem 1.2 states that no family of ten vectors in $\mathbb C^4$ possesses the phase-retrieval property \cite{Huang}.  Corollary 1.3 combines this lower bound with Vinzant's explicit eleven-vector frame \cite{Vinzant} and concludes
\[
 N_{\min}(\mathbb C^4)=11.
\]
The theorem and corollary appear on supporting PDF page 4; the lower-bound proof is on pages 10--14.  Vinzant's frame and its injectivity certificate appear on the second supporting PDF, pages 1--2.

\section*{Proof intuition}
Suppose ten vectors did phase retrieval.  After Parseval normalization, their ten intensities sum to one, so they define an injective smooth map from $\mathbb{CP}^3$ into a nine-dimensional affine hyperplane.  The rank-two kernel criterion for phase retrieval makes its differential injective; compactness then turns the map into a smooth embedding
\[
 G:\mathbb{CP}^3\hookrightarrow\mathbb R^9.
\]
Its oriented normal bundle $\nu$ has real rank three.  The stable tangent--normal identity
\[
 T\mathbb{CP}^3\oplus\nu\cong\underline{\mathbb R}^{,9}
\]
and the Euler sequence give $p_1(T\mathbb{CP}^3)=4h^2$, hence $p_1(\nu)=-4h^2$.

Now use the special rank-one form of an intensity coordinate.  For any nonzero measurement vector $a_j$, its zero locus is a projective hyperplane $H_j\cong\mathbb{CP}^2$.  Along $H_j$, the corresponding centered coordinate direction is everywhere orthogonal to the embedded tangent space, yielding a nowhere-zero section of $\nu|_{H_j}$.  Thus
\[
 \nu|_{H_j}\cong\underline{\mathbb R}\oplus\eta
\]
for an oriented rank-two bundle $\eta$.  Writing $\bar h=h|_{H_j}$, one side forces
\[
 p_1(\nu|_{H_j})=-4\bar h^2,
\]
while the splitting forces $p_1(\nu|_{H_j})=p_1(\eta)=e(\eta)^2=k^2\bar h^2$ for some integer $k$.  The impossible equation $k^2=-4$ rules out ten measurements.  Vinzant's certified frame supplies eleven, closing the gap.

\section*{Scope boundary}
This is an explicit later resolution, not an agent-generated theorem.  It settles the source's standard rank-one question only in complex dimension four.  It does not determine the separate minima for generalized Hermitian measurements, fusion-frame projections, arbitrary low-rank recovery, or complex dimensions $d\ge6$ not covered by the known $d=2^k+1$ family.

{\raggedright
\begin{thebibliography}{9}
\bibitem{WX}
Y. Wang and Z. Xu,
\emph{Generalized phase retrieval: measurement number, matrix recovery and beyond},
arXiv:1605.08034 (2016/2019), p.~5.

\bibitem{Huang}
M. Huang,
\emph{Phase Retrieval in $\mathbb C^4$ Requires Exactly Eleven Measurements},
arXiv:2607.27719v1 (2026), Theorem~1.2 and Corollary~1.3.

\bibitem{Vinzant}
C. Vinzant,
\emph{A small frame and a certificate of its injectivity},
arXiv:1502.04656 (2015), Theorem~1.
\end{thebibliography}
\par}

\end{document}
